\documentclass[11pt]{article}
\usepackage[margin=0.9in]{geometry}
\usepackage{array}
\usepackage{booktabs}
\usepackage{amsmath,amssymb}
\usepackage{graphicx}
\usepackage{caption}
\usepackage[hidelinks]{hyperref}
\newcommand{\code}[1]{\texttt{#1}}
\captionsetup{font=small,skip=4pt}

\title{SAC vs SAC+JEPA with a cosine-normalized one-step objective\\[2pt]
\large A 10{,}000-transition RTX~3090 comparison}
\author{}
\date{September 20, 2026}

\begin{document}
\maketitle

\noindent
\textbf{Summary.} Normalizing the predictive loss to a scale-free cosine
distance fixes the loss-scale defect of the earlier MSE branch and yields an
auxiliary objective that is demonstrably learning: at 10k transitions the
predictor reaches $0.0038$ against a shuffled-action loss of $0.0099$ and a
no-model identity baseline of $0.1253$. Performance, however, remains below
vanilla SAC --- $1.188$ versus $1.655$ --- in this single-seed, 10k-transition
comparison.

\section{Shared setup}

All three arms use identical terrain (standardized centered-Gaussian,
periodic tiling, random start positions), the $P_{\max}=6$ longitudinal power
cap, 16 parallel environments, the same replay configuration
($20{,}000$ capacity, batch $128$, $1{,}000$ learning-start transitions, one
update per transition), the learned entropy temperature, the same evaluation
interval of $2{,}048$ transitions, and held-out terrain seeds
$100001$--$100004$. Only the actor objective differs. Each arm is a single
training seed; the reported $\pm$ are standard deviations across the four
held-out terrain seeds, not across runs.

\section{Full objective}

The actor minimizes
\begin{equation}
 \mathcal L_{\rm total}
 =\underbrace{\mathbb E_{o\sim\mathcal D}\big[\alpha\log\pi_\theta(a|o)-\min(Q_1,Q_2)(o,a)\big]}_{\text{SAC actor loss}}
 +\lambda_{\rm JEPA}\,\mathcal L_{\rm JEPA}
 +c_{\rm var}\,\mathcal L_{\rm var},
 \qquad a\sim\pi_\theta(\cdot|o),
\end{equation}
with the scale-free predictive term
\begin{equation}
 \mathcal L_{\rm JEPA}
 =1-\cos\!\Big(q_\psi\big(\big[f_\theta(o_t),\,a_t\big]\big),\;
 \operatorname{sg}\big[\bar f(o_{t+1})\big]\Big),
\end{equation}
and the VICReg-style variance floor
\begin{equation}
 \mathcal L_{\rm var}
 =\frac{1}{d}\sum_{j=1}^{d}\operatorname{relu}\!\Big(c_{\rm var}^{\rm tgt}
 -\sqrt{\operatorname{Var}_{b}\big[f_\theta(o_t)\big]_j+\epsilon}\Big),
 \qquad \epsilon=10^{-4}.
\end{equation}

\noindent
Here $f_\theta$ is the online actor encoder, $q_\psi$ a $130\!\to\!128$ predictor
consuming the current embedding and the executed action, $\bar f$ a
stop-gradient EMA copy of $f_\theta$, and $\operatorname{sg}$ the stop-gradient.
SAC's target \emph{critics} use $\tau=0.005$; this EMA target is separate and
updates once per $64$ collected transitions with $\tau=0.01$. Replay supplies
the executed tanh-squashed action and the true pre-reset next observation.
Because $\cos$ discards magnitude, $\mathcal L_{\rm var}$ is what prevents
collapse; the two terms are complementary, not alternatives. Substituting MSE
for the cosine gives the ``corrected MSE'' arm.

\begin{center}
\begin{tabular}{lccc}
\toprule
\textbf{Setting} & \textbf{SAC} & \textbf{SAC+JEPA (MSE+var)} & \textbf{SAC+JEPA (cosine)} \\
\midrule
$\lambda_{\rm JEPA}$ & $0$ & $0.025$ & $0.025$ \\
predictive loss & --- & MSE on embeddings & $1-\cos$ \\
$c_{\rm var}$ / target & --- & $0.01$ / $0.1$ & $0.01$ / $0.1$ \\
EMA cadence & --- & per $64$ transitions & per $64$ transitions \\
\bottomrule
\end{tabular}
\end{center}

\section{Results}

\begin{center}
\begin{tabular}{lccc}
\toprule
\textbf{Transitions} & \textbf{SAC} & \textbf{SAC+JEPA (MSE+var)} & \textbf{SAC+JEPA (cosine)} \\
\midrule
2{,}048  & $0.064\pm0.031$ & $0.065\pm0.030$ & $0.064\pm0.029$ \\
4{,}096  & $0.190\pm0.027$ & $0.174\pm0.028$ & $0.275\pm0.035$ \\
6{,}144  & $0.292\pm0.042$ & $0.286\pm0.036$ & $0.269\pm0.027$ \\
8{,}192  & $0.673\pm0.041$ & $0.701\pm0.033$ & $0.651\pm0.030$ \\
10{,}000 & $\mathbf{1.655\pm0.056}$ & $\mathbf{1.403\pm0.020}$ & $\mathbf{1.188\pm0.042}$ \\
\bottomrule
\end{tabular}
\end{center}

\noindent
The endpoints are ordered monotonically in the influence of the auxiliary term:
no JEPA ($1.655$) $>$ MSE JEPA ($1.403$, $-15.2\%$) $>$ cosine JEPA
($1.188$, $-28.2\%$). All three arms share an effectively identical entropy
trajectory ($\alpha$: $0.1457,\,0.0788,\,0.0426,\,0.0231,\,0.0135$) and
near-identical critic loss, so the divergence is late and confined to the
policy representation.

\section{The objective is learning, not collapsing}

The cosine arm logs three quantities per checkpoint that separate a working
predictor from a degenerate one: the actual loss, the loss with actions
permuted across the batch, and the no-model identity baseline
$1-\cos(z_t,z_{t+1})$.

\begin{center}
\begin{tabular}{lcccc}
\toprule
\textbf{Transitions} & $\mathcal L_{\rm JEPA}$ & shuffled action & identity baseline & $\lambda\mathcal L_{\rm JEPA}$ \\
\midrule
2{,}048  & $0.00801$ & $0.01863$ & $0.26834$ & $2.0\times10^{-4}$ \\
4{,}096  & $0.00355$ & $0.01167$ & $0.35455$ & $8.9\times10^{-5}$ \\
6{,}144  & $0.00292$ & $0.01080$ & $0.19969$ & $7.3\times10^{-5}$ \\
8{,}192  & $0.00264$ & $0.00944$ & $0.15826$ & $6.6\times10^{-5}$ \\
10{,}000 & $0.00376$ & $0.00986$ & $0.12528$ & $9.4\times10^{-5}$ \\
\bottomrule
\end{tabular}
\end{center}

\noindent
At 10k the predictor is $33\times$ better than the identity baseline and
$2.6\times$ better than the action-shuffled variant, so it is neither
collapsed nor ignoring the action. The action-conditioning is nonetheless
secondary: the shuffled predictor already beats identity by $12.7\times$, so
most of the predictive power comes from $z_t$ alone.

Two caveats bound the interpretation. First, the weighted auxiliary
contribution $\lambda\mathcal L_{\rm JEPA}\approx9.4\times10^{-5}$ is
$0.003\%$ of the actor loss ($-3.304$); a loss this small relative to the
policy term can still carry a disproportionate \emph{gradient} because the
cosine gradient scales with $1/\lVert z\rVert$, so loss magnitude alone does
not establish that the term is inert. Second, the variance floor is active only
at the first checkpoint ($\mathcal L_{\rm var}=0.011$, then $0$ thereafter, with
embedding standard deviation $0.09$--$0.12$ above the $0.1$ target), so the
comparison is effectively SAC-only versus MSE-JEPA versus cosine-JEPA at
$\lambda=0.025$.

\section{Performance plots}

\begin{figure}[h]
\centering
\includegraphics[width=0.32\linewidth]{../../../dumps/sac_10k/plots/learning_curve.png}
\includegraphics[width=0.32\linewidth]{../../../dumps/sac_jepa_corrected_10k/plots/learning_curve.png}
\includegraphics[width=0.32\linewidth]{../../../dumps/sac_jepa_cosine_10k/plots/learning_curve.png}
\caption{Learning curves. Left: SAC. Centre: SAC+JEPA MSE+variance. Right:
SAC+JEPA cosine. Lower panels are the auxiliary objective.}
\end{figure}

\begin{figure}[h]
\centering
\includegraphics[width=0.48\linewidth]{../../../dumps/sac_jepa_cosine_10k/plots/dynamics_history_sac_000010000_seed100001.png}
\includegraphics[width=0.48\linewidth]{../../../dumps/sac_jepa_cosine_10k/plots/acceleration_history_sac_000010000_seed100001.png}
\caption{Cosine arm, held-out seed $100001$: physical and latent velocity with
deterministic controls, and acceleration magnitudes with signed
parallel/perpendicular components.}
\end{figure}

\begin{figure}[h]
\centering
\includegraphics[width=0.48\linewidth]{../../../dumps/sac_jepa_cosine_10k/plots/inference_trajectories_seed100001.png}
\includegraphics[width=0.48\linewidth]{../../../dumps/sac_jepa_cosine_10k/plots/inference_histories_seed100001.png}
\caption{Cosine arm: deterministic inference trajectories over the terrain and
the corresponding state/action histories.}
\end{figure}

\begin{figure}[h]
\centering
\includegraphics[width=0.48\linewidth]{../../../dumps/sac_jepa_cosine_10k/plots/episode_diagnostics_sac_000010000_seed100001.png}
\includegraphics[width=0.48\linewidth]{../../../dumps/sac_jepa_cosine_10k/plots/latent_umap_sac_000010000_seed100001.png}
\caption{Cosine arm: full episode diagnostic and endpoint actor-encoder UMAP
coloured by physical speed.}
\end{figure}

\begin{figure}[h]
\centering
\includegraphics[width=0.62\linewidth]{../../../dumps/sac_jepa_cosine_10k/plots/pairwise_encoder_umap_5x5_seed100001.png}
\caption{Cosine arm: pairwise joint UMAP grid across checkpoints $2{,}048$
through $10{,}000$ on one common held-out trajectory.}
\end{figure}

\begin{figure}[h]
\centering
\includegraphics[width=0.44\linewidth]{../../../dumps/sac_10k/plots/encoder_cka_matrix_seed100001.png}
\includegraphics[width=0.44\linewidth]{../../../dumps/sac_jepa_cosine_10k/plots/encoder_cka_matrix_seed100001.png}
\caption{Linear CKA between checkpoint encoders on a held-out trajectory.
Left: SAC. Right: SAC+JEPA cosine.}
\end{figure}

\clearpage
\section{Conclusion}

The loss-scale defect is fixed. Cosine normalization makes $\lambda$ a
scale-free weight, the adversarial diagnostics show a predictor that genuinely
uses both $z_t$ and $a_t$, and the objective is no longer the trivial
near-identity map diagnosed in the earlier MSE branch.

The auxiliary objective nevertheless still trails vanilla SAC at this budget:
$1.188$ against $1.655$ ($-28.2\%$), with the MSE variant intermediate at
$1.403$. Three limits prevent a stronger claim. The comparison is
single-seed, so with $\pm$ reported across held-out terrains rather than runs
the gaps are not statistically resolvable. The $10$k budget sits inside the
takeoff region, where a shift in takeoff timing moves the endpoint by roughly
the observed margin. And the horizon is still a single step, so the target
carries little long-range dynamics even though it is now predicted well.

Establishing whether one-step JEPA helps or harms SAC requires matched
multi-seed runs at $\ge100$k transitions, with the variance floor and the
action-conditioning ablated independently.

\end{document}
